\vspace{-15pt}
\section{Introduction}\label{I}
Autonomous target search is a cornerstone of modern drone operations, underpinning critical applications such as disaster response~\cite{schedl2021autonomous}, aerial surveillance~\cite{sun2024moving}, and environmental surveying~\cite{geng2025epic}. 
However, conventional search strategies have been constrained by closed-set detectors and geometry-driven heuristics~\cite{luo2024star,schedl2021autonomous}. The recent emergence of Vision-Language Models (VLMs) has catalyzed a paradigm shift toward open-vocabulary semantic object navigation. This evolution empowers drones to interpret natural language instructions and leverage environmental context, unlocking the potential to fundamentally enhance search efficiency in complex environments.
Despite this promise, integrating VLMs into existing aerial systems remains challenging in practice.

The first challenge stems from the system design of existing VLM-based search systems, which typically treat VLMs as step-wise planners to generate discrete actions~\cite{ji2025towards,xiao2025uav}. Although effective in confined settings, this action-output design introduces two fundamental bottlenecks. 
(\textit{i}) Frequency mismatch: VLMs' high inference latency (often over 2000ms~\cite{wang2025uav}) inherently limits them to sparse action outputs, creating a severe frequency mismatch with drones' real-time planning modules ($>$10Hz, 100ms). Current ``stop-and-infer'' paradigms force drones to hover for VLM responses, disrupting continuous flight, reducing efficiency, and raising mission failure risks in large-scale scenarios given the limited flight endurance.
(\textit{ii}) Limited 3D scene understanding: Since VLMs infer based on 2D visual observations, they often struggle to associate visual entities across viewpoints and integrate them into a global spatial context, resulting in spatially inconsistent actions as the viewpoint shifts. These unstable actions are exacerbated in complex and large-scale outdoor scenarios during long-horizon missions, resulting in premature abandonment of promising regions and redundant revisits.

Secondly, existing planning strategies lack a unified optimization framework that can dynamically balance VLM-derived semantic guidance with geometric motion efficiency. Efficient aerial object navigation requires both effective exploitation of semantic cues and rational consideration of motion costs, yet current methods mostly adopt greedy decisions or simplistic fusion schemes that fail to reconcile both aspects in diverse environments. This absence of systematic balancing induces a binary trade-off in practice: overemphasizing coarse and uncertain semantic guidance results in circuitous routes that exhaust battery resources, while favoring geometric efficiency leads to myopic exploration that overlooks critical regions. This limitation is particularly pronounced in expansive outdoor environments with sparse semantic cues and diverse layouts, severely restricting search performance.

\textbf{Our work}. We tackle the above challenges and propose \sysname, achieving smooth and efficient aerial object navigation in outdoor environments by seamlessly bridging VLM's semantic reasoning with continuous path planning. As shown in \fig\ref{fig:teaser}, \sysname can efficiently locate \textit{open-set} objects with \textit{zero-shot} generalization for various scenarios. \sysname features key designs in both system architecture and algorithms to fully unleash the VLM's potential, as elaborated below.


On the architectural front, we advocate a simple yet effective design philosophy: repositioning the VLM as a high-level dense semantic generator rather than a direct action generator, and integrating semantic information through a 3D representation. 
Since the direct action-generation paradigm with high-latency VLMs cannot meet the high-frequency requirements of real-time flight, we instead use the VLM to predict region-level target-presence probabilities from selected keyframes, producing dense semantic priors that provide more informative guidance for motion planning.
To mitigate limitations in 3D scene understanding, we fuse these language-conditioned priors into a 3D value map as a temporal integrator, which uses confidence-weighted, cross-view integration to reconcile potentially contradictory VLM outputs over time.
Meanwhile, this 3D value map establishes a principled interface that supports a \textit{dual-pathway asynchronous} design between VLM reasoning and path planning, enabling effective synergy while both modules run at their native frequencies and alleviating the latency-mismatch bottleneck.
Specifically, the VLM-driven reasoning pathway asynchronously embeds semantic understanding into the 3D value map, while the high-frequency planning pathway continuously harvests this evolving semantic memory to generate motion trajectories without interruption.


On the algorithm front, we complement the above architecture with two targeted optimizations. 
(\textit{i}) \textit{Active Dual-Task Reasoning (ADTR)}: while the architecture enables asynchronous inference, the continuous drone movement generates high-frequency image streams with redundant content that renders per-frame processing computationally prohibitive. ADTR thus selectively queries the VLM only when observations provide novel geometric coverage or contain instruction-relevant information. This adaptive mechanism turns the VLM into an active reasoner, substantially reducing unnecessary computation while maintaining scene awareness.
(\textit{ii}) \textit{Semantic-Geometric Coherent Planning (SGCP)}: we propose SGCP to resolve the inherent tension between semantic guidance and motion efficiency. SGCP introduces a selective constraint injection mechanism that activates precedence constraints only between regions whose semantic values differ significantly, then minimizes the total travel distance subject to these constraints. This establishes a dynamic partial order: high-value regions are prioritized when their semantic relevance is distinct; when values are comparable, visitation order is flexibly determined by geometric proximity. 
By framing semantic-geometric reconciliation as a unified constrained optimization under value-gap-derived precedence constraints, SGCP adapts seamlessly to environmental heterogeneity, exhibiting semantic-driven behavior in scenarios with distinct high-value regions and geometry-driven efficiency in homogeneous regions.


We conduct extensive experiments in a high-fidelity simulator Unreal Engine~\cite{ue} with diverse object navigation tasks across various environments such as urban downtowns and wilderness villages. 
Evaluation results show that \sysname achieves an average success rate of 73\%, and an average navigation error of 11.6m, outperforming baselines by 49.1\% and 80.3\%, respectively. \sysname also reduces the total flight time by 59.2\%. We further validate our system through over ten hours of real-world experiments, deploying \sysname on a customized quadrotor across diverse and challenging environments, thereby demonstrating its practical effectiveness.

In summary, this paper makes the following contributions.
\begin{itemize}
    \item We propose \sysname, an outdoor open-set aerial object navigation system that repositions VLMs as dense semantic generators rather than direct action generators, and integrates semantics via a 3D value map. It further features a dual-pathway asynchronous design that alleviates latency mismatch and VLMs’ limited 3D scene understanding, enabling continuous flight with adaptive semantic guidance.
    \item We propose ADTR that exploits geometric and semantic redundancy for selective VLM querying; and SGCP, a unified optimization framework that dynamically reconciles semantic priorities and motion efficiency, enabling seamless adaptation to heterogeneous environments.
    \item We evaluate \sysname in a large variety of navigation tasks across various types of environments to demonstrate its zero-shot generalization. We further deploy \sysname in challenging real-world scenarios with extensive evaluation to demonstrate its superior performance. 
\end{itemize}


